\documentclass[12pt]{memoir}

% Match the format of the ClassifyingStack summary.
\def\nsemestre{II}
\def\nterm{September}
\def\nyear{2026}
\def\nprofesor{}
\def\nauthor{Ignacio Rojas}
\def\nsigla{RR0915}
\def\nsiglahead{Riemann--Roch and very ampleness}
\def\nextra{Expository notes}
\def\nlang{ENG}
\let\footruleskip\relax
\input{../../../headerVarillyDiff.tex}
\author{Ignacio Rojas\\[0.4ex]ChatGPT (OpenAI)}

\newcommand{\RRtag}[1]{\href{https://stacks.math.columbia.edu/tag/#1}{Tag~\texttt{#1}}}
\newcommand{\RRcheckpoint}[1]{\par\smallskip{\raggedright\noindent\blu{\textbf{Checkpoint.}} #1\par}\smallskip}
\newcommand{\RRsolution}[1]{\subsection*{Solution to Exercise~\ref{#1}}\label{sol:#1}}
\counterwithout{section}{chapter}
\setsecnumdepth{subsection}
\settocdepth{section}
\setlength{\columnsep}{20pt}
\setlength{\emergencystretch}{1.5em}
\renewcommand{\bibsection}{\section{References}\label{sec:references}}
\hypersetup{pdftitle={Riemann--Roch and very ampleness: RR0915},
  pdfauthor={Ignacio Rojas; ChatGPT (OpenAI)},
  pdfsubject={Sections, negative degree, point and tangent separation, and pseudostability}}

\begin{document}
\maketitle
\begin{center}
\emph{From a degree bound to an embedding of a curve}
\par\bigskip
\begin{minipage}{0.8\textwidth}
\small\centering
Developed from Ignacio Rojas's questions following an advisor meeting,
with exposition and drafting by ChatGPT (OpenAI).
The follow-up question about negative degree is placed at the first
use of that argument. Exercises and separate solutions check both
the calculations and their geometric meaning.
\end{minipage}
\end{center}
\clearpage
\tableofcontents*
\clearpage
\begin{multicols}{2}
\raggedcolumns

\section{Context and the goal}\label{sec:context}

The original question concerned a numerical condition for
pseudostability of a marked curve. Write
\[
 M=p_1+\cdots+p_n,\qquad B=\omega_C(M).
\]
For connected reduced projective curves with only nodes and ordinary
cusps, distinct smooth markings, and $g+n>2$, our discussion related
pseudostability to very ampleness of $B^{\otimes2}$.
The relevant embedding theorem of Catanese--Franciosi--Hulek--Reid
(CFHR) checks degrees on \emph{all subcurves}.

During the advisor meeting, the smooth-curve version was explained
using Riemann--Roch and Serre duality. The purpose here is to fill in
the gaps in that proof: why negative degree rules out sections, why
dimensions detect base points and distinguish points, and why a
simple zero detects a tangent direction. Section~\ref{sec:pseudo}
returns to the original singular-curve problem.

\begin{Th}\label{th:embedding}
Let $C$ be a smooth connected projective curve over $\bC$ of genus $g$.
If $L$ is a line bundle with $\deg L\ge 2g+1$, then $L$ is very ample.
\end{Th}

Unless stated otherwise, these smoothness and projectivity assumptions
remain in force. Set $d=\deg L$ and
\[
 h^i(L)=\dim_{\bC}H^i(C,L).
\]
In particular, $h^0$ counts \emph{linearly independent sections}, not
the number of elements of the space of sections. We write
$L(-D)=L\otimes\cO_C(-D)$, and $K_C$ denotes a canonical divisor,
so $\cO_C(K_C)\cong\omega_C$.

\section{Degree, zeros, and the apparent contradiction}\label{sec:degree}

\subsection{What the order of a section means}\label{sec:local-sections}
First recall how to write a section using ordinary functions. A line
bundle has a one-dimensional vector space $A|_x$ over each point $x$.
Near any point, on a sufficiently small open set $U$, it can be
identified with a product:
\[
 A|_U\cong U\times\bC.
\]
This identification is called a \emph{local trivialization}. It
preserves the point of $U$ and is linear on each fiber. A vector in
the fiber over $x$ is now described by $(x,a)$, where $a\in\bC$.
A regular section on $U$ is consequently described by
\[
 x\longmapsto(x,f(x))
\]
for a regular function $f$ on $U$. Its value at $x$ is zero exactly
when $f(x)=0$. Thus, after choosing local coordinates in the fibers,
we can calculate with the function $f$ representing the section.

Choosing such an identification is equivalent to choosing a
nonzero vector, hence a \emph{basis}, in each one-dimensional fiber,
varying regularly with $x$. This is a choice of bases for the fibers,
not a basis of the vector space $H^0(C,A)$ of global sections.
The product description is local; a single such identification
need not exist over the whole curve.

Now let a nonzero regular section $s$ be represented locally by $f$.
In a local parameter $z$ vanishing at $x$, write
\[
 f=z^m u,\qquad u(x)\ne0.
\]
Then $\operatorname{ord}_x(s)=m\ge0$. A different local trivialization
multiplies the representing function by a nowhere-zero regular
function, which does not change its order of vanishing. Since $s$
is not identically zero, it has only finitely many zeros and defines
the effective divisor
\[
 D_s=\sum_{x\in C}\operatorname{ord}_x(s)\,x.
\]
Here $s$ is a \emph{section}; $D_s$ is its \emph{divisor of zeros}.
``Effective'' describes the nonnegative coefficients of $D_s$.

Multiplication by $s$ identifies $\cO_C(D_s)$ with $A$.
To see this in functions, shrink $U$ so that $u$ is nowhere zero.
Local sections of $\cO_C(D_s)$ have the form $z^{-m}h$, with $h$
regular. Multiplication by the function $f=z^m u$ representing $s$
sends
\[
 z^{-m}h\longmapsto uh.
\]
This is an isomorphism onto the functions representing sections of
$A$, since multiplication by $u$ has inverse multiplication by
$u^{-1}$. Consequently,
\[
\begin{aligned}
 A&\cong\cO_C(D_s),\\
 \deg A&=\sum_x\operatorname{ord}_x(s)\ge0.
\end{aligned}
\]

\begin{Lem}\label{lem:negative}
If $\deg A<0$, then $H^0(C,A)=\{0\}$.
\end{Lem}
\begin{proof}
If a nonzero regular section existed, its effective zero divisor
would give $\deg A\ge0$, a contradiction.
\end{proof}

\subsection{The follow-up question, at its point of use}
\begin{Qn}
If $\deg A$ is a sum of nonnegative orders, how could that sum equal
a negative degree?
\end{Qn}

It cannot. The statement about a nonnegative sum is \emph{conditional
on the existence of a nonzero regular section}. For example, if
$\deg A=-1$, assuming such an $s$ exists would force
\[
 -1=\sum_x\operatorname{ord}_x(s)\ge0.
\]
This impossibility disproves the existence of $s$. It does not
construct a negative sum of nonnegative numbers.

The degree of $A$ is defined even when $A$ has no nonzero regular
sections. One way to compute it is to choose a nonzero meromorphic
section $t$ and form its divisor. Then
\[
 \deg A=\sum_x\operatorname{ord}_x(t),
\]
where poles contribute negative orders. Equivalently, degree is the
total zero multiplicity minus the total pole multiplicity. Different
choices give the same degree: their ratio is a meromorphic function,
whose principal divisor has degree zero.

Thus a negative-degree line bundle exists perfectly well; every
nonzero meromorphic section has more poles than zeros, and none can
be regular everywhere. The zero section always exists, but has no
finite divisor of zeros to which this argument applies.

\RRcheckpoint{State the argument in the order ``negative degree;
assume a nonzero regular section; obtain a contradiction.'' Which
assumption is rejected?}

\section{The vanishing and dimension calculation}\label{sec:rr}

For a line bundle $F$, Serre duality and Riemann--Roch give
\[
\begin{aligned}
 H^1(C,F)&\cong H^0(C,\omega_C\otimes F^{-1})^\vee,\\
 h^0(F)-h^1(F)&=\deg F+1-g.
\end{aligned}
\]
The first formula changes an $H^1$ computation into a question about
sections. Lemma~\ref{lem:negative} then answers that question by degree.
The second formula computes $h^0$ once $h^1$ vanishes.

For $F=L$, the degree of the dual bundle is
\[
 \deg(\omega_C\otimes L^{-1})=2g-2-d\le-3.
\]
Hence $H^1(C,L)=0$. The calculation
$(2g-2)-(2g-1)=-1$ from the meeting uses the weaker bound
$d\ge2g-1$, which suffices for this first vanishing. The stronger
bound $d\ge2g+1$ allows us to subtract two points as well.

\begin{Lem}\label{lem:uniform}
Let $d\ge2g+1$ and let $D$ be an effective divisor of degree $r\le2$.
Then
\[
\begin{aligned}
 H^1(C,L(-D))&=0,\\
 h^0(L(-D))&=d-r+1-g.
\end{aligned}
\]
\end{Lem}
\begin{proof}
The line bundle dual to $L(-D)$ in Serre duality has degree
\[
\begin{aligned}
 \deg(\omega_C\otimes L^{-1}(D))
 &=2g-2-d+r\\
 &\le r-3<0.
\end{aligned}
\]
It has no nonzero sections. Duality gives the vanishing, and
Riemann--Roch gives the claimed dimension.
\end{proof}

In particular, for distinct $p,q\in C$,
\begin{equation}\label{eq:dimensions}
\begin{aligned}
 h^0(L)&=d+1-g,\\
 h^0(L(-p))&=d-g,\\
 h^0(L(-p-q))&=d-1-g,\\
 h^0(L(-2p))&=d-1-g.
\end{aligned}
\end{equation}
Each additional required order of vanishing removes one independent
section. The next section explains why these drops give an embedding.

\section{The three geometric tests}\label{sec:geometry}

\subsection{Base points: can a section be nonzero at $p$?}\label{sec:basepoints}
We want to show that, for every $p$, some global section of $L$ is
nonzero at $p$. The comparison with $L(-p)$ uses two different maps:
an inclusion of sheaves and an evaluation map on global sections.
The inclusion exists for every $L$, independently of any assumption
about the dimensions of its spaces of sections.

\subsubsection*{The inclusion on sheaves and the map on fibers}
We use the same notation $L$ for a line bundle and its sheaf of local
sections. The arrow $L(-p)\hookrightarrow L$ is an injection of these
\emph{sheaves}; it need not be injective on every fiber.

On a smooth curve, $\cO_C(-p)$ is the ideal sheaf of the point $p$:
its sections on a neighborhood of $p$ are regular functions vanishing
at $p$. Its natural inclusion into $\cO_C$, tensored with the locally
free sheaf $L$, gives
\[
 L(-p)=L\otimes\cO_C(-p)\hookrightarrow L.
\]
This is the ideal-sheaf interpretation of a negative effective
divisor; see \RRtag{0C4T} in \cite{Stacks}.

Use the local description from Section~\ref{sec:local-sections}:
identify $L|_U$ with $U\times\bC$, so a section is represented by
a regular function $f$. Choose a local parameter $z$ with $z(p)=0$.
The section vanishes at $p$ precisely when $f(p)=0$, in which case
$f=zg$ for a regular $g$, after shrinking $U$ if needed.

This gives a compatible local description of $L(-p)$ by functions
$g$: when its section is viewed as a section of $L$, the representing
function is $zg$. Thus the map on local sections is simply
\[
 g\longmapsto zg.
\]
The input function describes a section of $L(-p)$; the output
function describes its image in $L$. Both bundles look like
$U\times\bC$ locally, but the map between them is multiplication
by $z$, not the identity.

It is injective on local sections: if $zg$ is \emph{identically}
zero, then $g$ is identically zero. In these product coordinates,
the corresponding bundle map is
\[
 \iota:(x,a)\longmapsto(x,z(x)a).
\]
On the fiber over $p$, this sends every $(p,a)$ to $(p,0)$ because
$z(p)=0$. In particular, the nonzero source vector $(p,1)$ has zero
image. After shrinking $U$ so that $z$ has no other zeros, the map
is an isomorphism on fibers away from $p$.

For example, the source section represented by the constant
function $1$ maps to the target section represented by $z$.
The latter is not the zero section: $z$ is not identically zero,
although $z(p)=0$. There is no contradiction between these assertions.
An injection of sheaves of sections is therefore different from a
fiberwise inclusion of vector bundles.

On global sections, $\iota$ induces an injection
\[
 H^0(C,L(-p))\hookrightarrow H^0(C,L).
\]
It takes sections of $L(-p)$ to sections of $L$ vanishing at $p$.
It does not take arbitrary sections of $L$ and force them to vanish.
Identifying the source with its image, we obtain
\[
\begin{aligned}
 &H^0(C,L(-p))\\
 &\qquad=\{s\in H^0(C,L):s(p)=0\}.
\end{aligned}
\]

\subsubsection*{The kernel belongs to the domain}
Give the three vector spaces different names:
\[
\begin{aligned}
 V&=H^0(C,L),\\
 W&=H^0(C,L(-p))\subseteq V,\\
 F&=L|_p.
\end{aligned}
\]
Here the inclusion of $W$ uses the identification just explained.
\begin{center}
\begin{tabularx}{\linewidth}{@{}lX@{}}
\toprule
Space & Meaning of an element\\
\midrule
$V$ & A global section of $L$.\\
$W$ & A global section of $L$ whose value at $p$ is zero.\\
$F$ & A single vector in the fiber over $p$.\\
\bottomrule
\end{tabularx}
\end{center}

Evaluation remembers just one value of a section:
\[
 \operatorname{ev}_p:V\longrightarrow F,\qquad s\longmapsto s(p).
\]
Its kernel is $W$, a subspace of the \emph{domain} $V$, not of the
target fiber $F$. Every section in $W$ maps to the same zero vector
of $F$, but the section itself can be nonzero elsewhere on $C$.
The arrangement of the spaces is
\[
 0\longrightarrow W\hookrightarrow V
 \xrightarrow{\operatorname{ev}_p}F.
\]
There is no terminal zero here until surjectivity has been proved.

Since $\dim F=1$, the \emph{image} of evaluation is either $\{0\}$
or all of $F$. Rank--nullity gives
\[
\begin{aligned}
 \operatorname{codim}_V W
 &=\dim V-\dim W\\
 &=\dim\operatorname{im}(\operatorname{ev}_p)\in\{0,1\}.
\end{aligned}
\]
Thus there are two possibilities:
\begin{itemize}
\raggedright
\item Codimension zero: $W=V$. Every section is in the kernel,
so every section vanishes at $p$. The point $p$ is a base point.
\item Codimension one: $\dim W=\dim V-1$. There is a section outside
the kernel, hence a section nonzero at $p$. The point $p$ is not a
base point.
\end{itemize}
Codimension one does not mean that $W$ is the zero space. If $V$ has
dimension $10$, such a kernel has dimension $9$.

\begin{Ex}[Many sections can have the same value zero]\label{ex:quadratic-evaluation}
Take $C=\bP^1$, $L=\cO_{\bP^1}(2)$, and $p=[0:1]$.
A global section is a homogeneous quadratic polynomial
\[
 s=aX^2+bXY+cY^2.
\]
On the chart $Y\ne0$, use $z=X/Y$. The local product description
of $\cO_{\bP^1}(2)$ represents this section by the regular function
obtained by dividing its homogeneous expression by $Y^2$:
\[
 \frac{s}{Y^2}=az^2+bz+c.
\]
The point $p$ has $z=0$, so evaluation gives the scalar $c$ in the
identified fiber $L|_p\cong\bC$. Consequently,
\[
\begin{aligned}
 V&=\operatorname{span}\{X^2,XY,Y^2\},\\
 W&=\operatorname{span}\{X^2,XY\}.
\end{aligned}
\]
The kernel $W$ has dimension two inside the three-dimensional space
$V$; it cannot be a subspace of the one-dimensional fiber $F$.
Both $X^2$ and $XY$ are nonzero sections whose values at $p$ are zero.
The section $Y^2$ has nonzero value there, so $p$ is not a base point.
\end{Ex}

\RRcheckpoint{For the local map $g\mapsto zg$, explain why it is
injective on sections but zero on the fiber at $p$. What happens
to the section represented by $g=1$? In the
evaluation map, which space contains the kernel, and which contains
the image?}

\subsubsection*{Using Riemann--Roch to finish the test}
There is no need to begin by assuming equality of dimensions.
Equation~\eqref{eq:dimensions} directly gives
\[
\begin{aligned}
 \dim V&=h^0(L)=d+1-g,\\
 \dim W&=h^0(L(-p))=d-g.
\end{aligned}
\]
Thus $W$ has codimension one in $V$: some section of $L$ is nonzero
at $p$. Since the same calculation holds for every $p$, $L$ is
basepoint-free.

\begin{Rmk}[Correcting the equality in the notes]
Basepoint-freeness requires a dimension drop, not equality of the
two dimensions. The sections nonzero at $p$ are the complement of
the kernel; they do not form a vector subspace.
\end{Rmk}

The sheaf exact sequence also explains the role of cohomology.
Writing $i:\{p\}\hookrightarrow C$, it is
\[
 0\to L(-p)\to L\to i_*(L|_p)\to0,
\]
where $i_*(L|_p)$ is the sheaf supported at $p$ with value $L|_p$.
The long exact sequence of cohomology contains
\[
 H^0(C,L)\longrightarrow L|_p
 \longrightarrow H^1(C,L(-p))=0.
\]
Thus every prescribed value at $p$ comes from a global section.

A basis $s_0,\ldots,s_N$ now defines a morphism
\[
\begin{aligned}
 \varphi_L:C&\longrightarrow\bP^N,\\
 x&\longmapsto[s_0(x):\cdots:s_N(x)].
\end{aligned}
\]
Basepoint-freeness ensures that these coordinates never all vanish.
A local trivialization turns the fiber values into numbers; changing
it rescales all coordinates together and leaves the projective
point unchanged.

\subsection{Distinct points: can a section distinguish $p$ and $q$?}
For $p\ne q$, there is an inclusion
\[
 H^0(C,L(-p-q))\subset H^0(C,L(-p)).
\]
The first space requires vanishing at both points; the second only
requires vanishing at $p$. Their dimensions differ by one, so we can
choose a section $s$ in the second space but outside the first. Then
\[
 s(p)=0,\qquad s(q)\ne0.
\]
Include $s$ in a basis. Its projective coordinate is zero at
$\varphi_L(p)$ and nonzero at $\varphi_L(q)$, so these images differ.
Equivalently, the hyperplane defined by $s$ contains one image but
misses the other. This proves injectivity of $\varphi_L$ on points.

\subsection{Tangents: can a section detect first-order movement?}
The inclusion
\[
 H^0(C,L(-2p))\subset H^0(C,L(-p))
\]
compares sections vanishing to order at least two with those
vanishing to order at least one. Again the dimension drops by one.
Hence there is an $s$ vanishing to order exactly one at $p$.

Choose $t$ with $t(p)\ne0$, using basepoint-freeness, and restrict
to a neighborhood where $t$ never vanishes. At each point $x$ there,
$s(x)$ is a unique scalar multiple of $t(x)$; those scalars form
the regular function $s/t$. Equivalently, represent the two sections
by functions in the same local trivialization and take their ratio.
In a local parameter $z$ centered at $p$,
\[
 \frac{s}{t}=az+O(z^2),\qquad a\ne0.
\]
The independent sections $s,t$ can be included in a basis, making
$s/t$ an affine coordinate of $\varphi_L$ near $p$. Its derivative
at $p$ is nonzero. Since $T_pC$ has dimension one, the differential
$d\varphi_{L,p}$ is injective.

This last check is necessary: the injective map
$z\mapsto(z^2,z^3)$ has zero derivative at the origin and produces a
cusp in its image. A simple zero supplies the nonzero linear term
that detects the tangent direction.

The complete linear system therefore has no base points and separates
points and tangent vectors. The projective embedding criterion gives
a closed embedding, proving Theorem~\ref{th:embedding}; see also
\RRtag{0H2V} in \cite{Stacks}.

\RRcheckpoint{Explain each strict inclusion by completing the sentence
``There is a section that vanishes \ldots\ but does not vanish
\ldots.'' For the tangent test, specify the order of vanishing.}

\section{Returning to pseudostability}\label{sec:pseudo}

Now allow the nodal and cuspidal curves from Section~\ref{sec:context}.
For a connected reduced union of components $Z\subseteq C$, put
\[
 h(Z)=p_a(Z),\qquad \ell(Z)=k(Z)+n(Z),
\]
where $k(Z)$ counts nodes attaching $Z$ to its complement and $n(Z)$
counts markings on $Z$. Internal singularities do not contribute to
$\ell(Z)$; their contribution is already reflected in $p_a(Z)$.
For $Z=C$, set $k(C)=0$.

Restriction of the dualizing bundle gives
\[
 \deg(B|_Z)=2h(Z)-2+\ell(Z).
\]
Under the stated assumptions, the usual conditions that $B$ be ample
and that there be no unmarked arithmetic-genus-one tails are
equivalently packaged as
\[
\begin{gathered}
 h(Z)+\ell(Z)>2\\
 \Longleftrightarrow\quad\deg(B|_Z)>h(Z)
\end{gathered}
\]
for every connected subcurve $Z$. For genus zero this requires at
least three attachments or markings; for genus one it requires at
least two. For proper subcurves of genus at least two it is automatic;
for $Z=C$ the assumption $g+n>2$ supplies it. Conversely the inequality
on each irreducible component gives positive degree, hence ampleness.

\subsection{Why the smooth theorem needs an extension}
A reducible curve can have a line bundle of very large total degree
whose restriction to one component is trivial. Such a component
cannot be embedded. Also, a section can vanish identically on one
component while being nonzero on another, so the smooth-curve
argument about its effective zero divisor cannot be used unchanged.

For our reduced curves, the sufficient part of CFHR's Theorem~1.1
\cite{CFHR} says that $A$ is very ample if
\[
 \deg(A|_Z)\ge2p_a(Z)+1
\]
for every connected subcurve $Z$. At singular points one tests all
length-two subschemes. On a smooth curve those tests are exactly
$p+q$ with $p\ne q$, and $2p$; at a singular point they require
more care because the point need not be a Cartier divisor.

For $A=B^{\otimes2}$, the integer inequality $\deg(B|_Z)>p_a(Z)$ gives
\[
 \deg(A|_Z)\ge2p_a(Z)+2.
\]
Thus CFHR proves the desired very ampleness for pseudostable curves.

Conversely, if $B^{\otimes2}$ is very ample, then $B$ is ample.
An unmarked genus-one tail $E$ attached at $q$ would have
$\deg(B|_E)=1$. Positivity on every irreducible component forces
$E$ to be irreducible. Its dualizing bundle is trivial, and
Riemann--Roch and Serre duality on this integral Gorenstein curve give
\[
\begin{aligned}
 B^{\otimes2}|_E&\cong\cO_E(2q),\\
 h^0(E,\cO_E(2q))&=2.
\end{aligned}
\]
This complete linear system maps to $\bP^1$ and cannot embed a
genus-one curve. Since a very ample bundle restricts to a very ample
bundle on a subcurve, such a tail is excluded.

The elementary proof explains the degree bound for smooth curves;
the subcurve theorem supplies the extra input for the original
pseudostability criterion. Merely requiring $B^{\otimes2}$ to be
\emph{ample} would add nothing to ampleness of $B$.

\section{Exercises: checking comprehension}\label{sec:exercises}

Attempt these before consulting Section~\ref{sec:solutions}.

\begin{Ej}[Negative degree and the missing assumption]\label{ex:negative}
On $\bP^1$, fix $p$ and let $A=\cO_{\bP^1}(-p)$.
What is $\deg A$? Explain why $H^0(\bP^1,A)=0$ in two ways:
using the contradiction argument, and by viewing a section as a
global regular function required to vanish at $p$.
Why does this not rule out a nonzero meromorphic section of $A$?
\end{Ej}

\begin{Ej}[Three degree thresholds]\label{ex:thresholds}
Use duality to show that $\deg F\ge2g-1$ implies $H^1(C,F)=0$.
Deduce that $\deg L\ge2g$ guarantees basepoint-freeness, and explain
why the argument for very ampleness uses $\deg L\ge2g+1$.
For $g=2$ and $d=5$, compute all four dimensions in
Equation~\eqref{eq:dimensions}.
\end{Ej}

\begin{Ej}[Evaluation and dimensions]\label{ex:evaluation}
Suppose $h^0(L)=4$. If $h^0(L(-p))=4$, what happens at $p$?
If instead $h^0(L(-p))=3$, what happens? Can this dimension be $2$?
Explain using the evaluation map and the dimension of its target.
\end{Ej}

\begin{Ej}[Separating two points]\label{ex:points}
Let $p\ne q$ and assume
$h^0(L(-p))>h^0(L(-p-q))$.
State exactly what section this produces. Assuming $L$ is
basepoint-free, explain why $\varphi_L(p)\ne\varphi_L(q)$.
On $\bP^1$ with coordinates $[X:Y]$ and $L=\cO(1)$, find such
a section for $p=[0:1]$ and $q=[1:0]$.
\end{Ej}

\begin{Ej}[A simple zero and a derivative]\label{ex:tangents}
If $s$ vanishes to order exactly one at $p$ and $t(p)\ne0$, compute
the first-order term of $s/t$ in a local parameter. What changes
if $s$ vanishes to order at least two? Explain why the map
$z\mapsto(z^2,z^3)$ is injective but has zero differential at $0$.
\end{Ej}

\begin{Ej}[Why total degree is insufficient]\label{ex:subcurves}
Let $C=C_1\cup R$, where both components are copies of $\bP^1$
meeting at one node. Glue $A|_{C_1}\cong\cO_{\bP^1}(m)$ and
$A|_R\cong\cO_R$, with $m\ge1$, by identifying their fibers at
the node. Compute $p_a(C)$ and $\deg A$. Why is $A$ not very ample,
even though the smooth numerical bound on total degree holds?
Which CFHR inequality fails?
\end{Ej}

\end{multicols}
\clearpage
\begin{multicols}{2}
\raggedcolumns
\section{Hints and solutions}\label{sec:solutions}

\RRsolution{ex:negative}
\hint{A regular function on a connected projective curve is constant.}
The degree is $-1$. A nonzero regular section would give an effective
zero divisor of degree $-1$, which is impossible. Alternatively,
$\cO(-p)\hookrightarrow\cO$ identifies its global sections with
global regular functions vanishing at $p$. The only such constant is
zero. Meromorphic sections may have poles, whose negative orders
allow their divisor to have degree $-1$.

\RRsolution{ex:thresholds}
\hint{After subtracting $r$ points, retain degree at least $2g-1$.}
We have
$\deg(\omega_C\otimes F^{-1})\le-1$, so duality gives $H^1(C,F)=0$.
For $d\ge2g$, this applies to both $L$ and $L(-p)$, and
Riemann--Roch gives the one-dimensional drop proving
basepoint-freeness. For $d\ge2g+1$, it also applies after subtracting
any effective divisor of degree two, giving both separation tests.
For $g=2,d=5$, the dimensions are $4,3,2,2$.
These are sufficient uniform bounds, not necessary conditions for
an individual line bundle.

\RRsolution{ex:evaluation}
\hint{Use rank--nullity for evaluation into $L|_p$.}
If the kernel has dimension $4$, evaluation is zero and $p$ is a
base point. If it has dimension $3$, evaluation is surjective and
some section is nonzero at $p$. Kernel dimension $2$ would require
rank $2$, impossible for a map to a one-dimensional space.

\RRsolution{ex:points}
\hint{Choose a vector in the larger space outside the smaller one.}
There is an $s$ with $s(p)=0$ and $s(q)\ne0$. Include it in a basis:
the corresponding coordinate vanishes at one projective image and
not at the other, so the images differ. In the stated example,
$s=X$ works. The basis $X,Y$ gives the identity embedding of
$\bP^1$ into $\bP^1$.

\RRsolution{ex:tangents}
\hint{Represent both sections by functions using the same local
product description of the line bundle.}
Let $s$ and $t$ be represented by $f=az+O(z^2)$ and $g=b+O(z)$,
where $a,b\ne0$. Then $s/t$ is the function
$f/g=(a/b)z+O(z^2)$, which has nonzero derivative. If $s$ has order at
least two, the ratio has no linear term. The map $(z^2,z^3)$ has
derivative $(2z,3z^2)$, zero at $0$. It is injective: away from the
origin the ratio of the second coordinate to the first recovers
$z$, and only $z=0$ maps to the origin.

\RRsolution{ex:subcurves}
\hint{A very ample line bundle remains very ample on a subcurve.}
The single-node union has arithmetic genus
$0+0+1-2+1=0$, while $\deg A=m\ge1=2p_a(C)+1$.
However, $A|_R\cong\cO_R$ has only constant sections and cannot
embed $R$. Hence $A$ is not very ample. The CFHR inequality on
$R$ would require
\[
 \deg(A|_R)=0\ge2p_a(R)+1=1,
\]
which fails. This example isolates the need to check subcurves.

\begin{thebibliography}{9}
\raggedright
\bibitem{Stacks}
The Stacks Project Authors, \emph{The Stacks Project}.
\href{https://stacks.math.columbia.edu/tag/0H2V}{Lemma~53.22.5,
Tag~\texttt{0H2V}}: very ampleness from degree at least $2g+1$,
using separation of length-two subschemes.
\RRtag{0C4T} identifies $\cO_C(-D)$ with the ideal sheaf of an
effective Cartier divisor $D$.

\bibitem{CFHR}
F.~Catanese, M.~Franciosi, K.~Hulek, and M.~Reid,
\href{https://doi.org/10.1017/S0027763000025381}{\emph{Embeddings of
curves and surfaces}}, Nagoya Mathematical Journal~154 (1999),
185--220. Theorem~1.1 gives the subcurve embedding criterion;
Remark~2.1 explains the length-two test.
\end{thebibliography}

\end{multicols}
\end{document}
